% =====================================================================
%  ch12.tex —— 第 12 章 相互作用费米子可逆相与自旋配边（原书 p.170–188）
%  §12.1 eta 不变量   §12.2 可逆相的一般低能有效作用量
%  §12.3 配边群       §12.4 配边与一般相互作用相
%  Part II 体例：无编号陈述环境，只有 \para 行内小标题与公式 (12.1)–(12.100)
%  表 10（p.183）与表 11（p.188）用 tabular 重排
% =====================================================================
\chapter{相互作用费米子可逆相与自旋配边}

% [原文 p.170]
% 注：p.170 上半页为第 11 章末尾（归 ch11 处理），本章从该页中部起。
我们已经看到，有能隙的自由费米子哈密顿量的分类给出了拓扑绝缘体与超导体的周期表，汇总在表 9 中。我们来看对角线上的 $\Z$，特别是 $i=d=1$（BDI 或 pin$^-$）、$i=d=2$（D 或 spin）与 $i=d=3$（DIII 或 pin$^+$）处的 $\Z$。

允许相互作用之后，$i=d=2$ 处的 $\Z$ 已知保持不变，而 $i=d=1$ 处的 $\Z$ 已知约化为 $\Z_8$，$i=d=3$ 处的 $\Z$ 已知约化为 $\Z_{16}$。$i=d=1$ 处的这一约化最早在 [FK09] 中被观察到，而我们在 §10.2.3 其实已经见过。的确，若体是 $1+1$ 维的，边界就只是一个 $0+1$ 维系统，那不过是量子力学系统的另一个名字。pin$^-$ 条件意味着 $T^2=+1$，即标准的时间反演对称性。当时我们看到，用自由（即二次的）哈密顿量无法使系统有能隙；但用八个 Majorana 费米子与四次哈密顿量，就可以让系统有能隙。

那么，相互作用费米子可逆相的一般分类理论是什么？如今普遍接受的答案是：它由配边群的某种对偶版本来给出，而这同样是代数拓扑的一大主题。本章的目的就是解释这些事情。

\section{eta 不变量}

任意的相互作用费米子系统很难概念化。但任何量子力学系统都应当在闭时空上有配分函数。我们就要用这样的配分函数来
% [原文 p.171]
刻画相互作用费米子可逆相。为此，需要先弄清质量极大的自由费米子系统的配分函数。

\subsection{定义}

取一个 $d+1$ 维的自由费米子系统，作 Wick 旋转，再把它放到流形 $M_{d+1}$ 上。它由如下形式的二次拉格朗日量给出：
\[
  S=\int d^{d+1}x\,\bar\psi(m+iD)\psi ,
\]
其中 $D$ 是 Dirac 算子。于是配分函数为
\begin{equation}
  Z(M_{d+1},m)=\int[d\psi][d\bar\psi]\,
  e^{-\int d^{d+1}x\,\bar\psi(m+iD)\psi}
  =\prod_a (m+i\lambda_a),
  \tag{12.1}
\end{equation}
其中 $\lambda_a$ 是 Dirac 算子 $D$ 的本征值。

我们只关心极端长距离极限，也就是取 $|m|\to\infty$。现在考虑比值
\begin{equation}
  \frac{Z(M_{d+1},m\to+\infty)}{Z(M_{d+1},m\to-\infty)}
  =\lim_{M\to+\infty}\prod_a\frac{i\lambda_a/M+1}{i\lambda_a/M-1},
  \tag{12.2}
\end{equation}
它比分子或分母单独一个都更好定义。

这是一个纯相位，因此把它写成
\begin{equation}
  \frac{1}{2\pi}\operatorname{Arg}
  \frac{Z(M_{d+1},m\to+\infty)}{Z(M_{d+1},m\to-\infty)}
  =\lim_{M\to+\infty}\sum_a\frac{1}{2\pi}\operatorname{Arg}
  \frac{i\lambda_a/M+1}{i\lambda_a/M-1},
  \tag{12.3}
\end{equation}
它在模 $1$ 意义下良定义。换一种提示性的写法，先对各个 $\lambda_a$ 取 $M\to\infty$ 的极限，则有
\begin{equation}
  =\frac{1}{2}\text{“}\sum_a\sgn\lambda_a\text{”}.
  \tag{12.4}
\end{equation}
当然，最后这个表达式需要正规化；但多亏重整化的一般理论，任何合理的正规化都行。

这称为 $M_{d+1}$ 上 Dirac 算子 $D$ 的 eta 不变量，记作 $\eta(D)$。等价地，
\begin{equation}
  Z(M_{d+1})=e^{2\pi i\eta(D)},
  \tag{12.5}
\end{equation}
其中左端是质量 $m\to+\infty$ 与 $m\to-\infty$ 的自由费米子系统配分函数之比。

几点评论：

\begin{itemize}
  \item eta 不变量过去也许吓到过你，但它不过是质量无穷大的自由费米子系统配分函数的相位部分，没什么可担心的。
% [原文 p.172]
  \item 在本讲义中，我们通常关心模 $1$ 的 eta 不变量，但它也可以定义为一个实数。这时按约定，$\lambda_a=0$ 的贡献定义为 $+1/2$。把它改成 $-1/2$ 不影响模 $1$ 的 eta 不变量，却会使 eta 不变量本身改变 $1$。
  \item 若去查数学原始文献，会发现 $\eta$ 的定义前面没有 $1/2$，零本征值的贡献为 $0$；这里写的东西在数学文献里其实记作 $\xi$。不过用 $\eta$ 表示 $\xi$ 如今在物理文献中已是标准做法，本讲义也照办。
  \item 数学文献通常用 zeta 函数正规化。但如前所述，任何合理的正规化都行；后面我们通常用 $e^{-s|\lambda_a|}$、$s\to0$ 的正规化。
\end{itemize}

\subsection{例：$S^1$ 上的复费米子}

作为例子，考虑 $S^1$ 上的复费米子，其 Dirac 算子为 $D=i\partial_\theta+A$，其中 $A$ 是常值的 $U(1)$ 规范场。取周期为 $\theta\sim\theta+2\pi$，则 $D$ 的本征值为 $\lambda_n=n+(a-1/2)$，$n\in\Z$，其中
\begin{equation}
  a-\frac{1}{2}:=\int_0^{2\pi}A\,d\theta .
  \tag{12.6}
\end{equation}
这里减去 $1/2$ 是一个约定，动机如下。考虑一个费米子系统并作 Wick 旋转，$\tr e^{-\beta H}$ 对应热时间方向上反周期的费米子；而周期的费米子对应的则是 $\tr(-1)^F e^{-\beta H}$。我们希望 $a=0$ 对应标准的热配分函数，而它应当是正的。

还要注意，谱 $\{n+(a-1/2)\}$ 在大规范变换 $a\to a+1$ 下不变。因此 eta 不变量也应当在 $a$ 的这一改变下不变。我们用这个自由度来保证 $|a|<1/2$。

于是，在谱中，$n+a-1/2$ 对 $n\ge1$ 为正，$n+a-1/2$ 对 $n\le0$ 为负。由此得到
\begin{equation}
  \sum_{n\ge1}(+1)e^{-s(n+1/2-a)}=\frac{1}{s}-a+O(s),\qquad
  \sum_{n\le0}(-1)e^{-s|n+1/2-a|}=\frac{1}{s}+a+O(s).
  \tag{12.7}
\end{equation}
所以
\begin{equation}
  \eta(D)=\lim_{s\to0}\frac{1}{2}\bigl(-2a+O(s)\bigr)=-a .
  \tag{12.8}
\end{equation}
$a=0$ 时 $\eta(D)=0$，这与上面的分析一致：它应当是热配分函数的相位，而热配分函数应当为实的正数。

注意 $\eta(D)$ 在 $a=\pm1/2$ 处不连续（由于大规范变换 $a\sim a+1$，这两个点应看作同一点），在那里它在 $\eta(D)=\pm1/2$ 之间跳变 $\pm1$。同时我们也发现，$e^{2\pi i\eta(D)}$、等价地 $\eta(D)$ 模 $1$，是连续的。

% [原文 p.173]
利用式 (12.6)，我们可以把关于 $\eta(D)$ 的结果写成
\begin{equation}
  \eta(D)\text{“}=\text{”}\frac{1}{2}-\int_{S^1}A\,d\theta .
  \tag{12.9}
\end{equation}
这里右端的第二项可以看作一维的 Chern--Simons 项。由于 $A$ 在大规范变换下不是规范不变的，右端看起来初等，其实并没有真正定义好，而且依赖规范。左端即 eta 不变量，则是用微分算子的谱与一个正规化的无穷和、以复杂得多的方式定义的；但按定义它显然是规范不变的。在这个意义上，可以把 eta 不变量看作 Chern--Simons 项的更好定义的版本。

\subsection{Atiyah--Patodi--Singer 指标定理：一个例子}

我们引入了流形 $M_{d+1}$ 上 Dirac 算子 $D_M$ 的 eta 不变量 $\eta(D_M)$ 这一概念。当 $M_{d+1}$ 是 $d+2$ 维流形 $W_{d+2}$ 的边界、即 $M_{d+1}=\partial W_{d+2}$ 时，它满足如下关系：
\begin{equation}
  \eta(D_M)=(\operatorname{index}D_N)+\int_N(\text{微分形式}),
  \tag{12.10}
\end{equation}
其中 $D_N$ 是 $N_{d+2}$ 上的 Dirac 算子，$\operatorname{index}D_N$ 是它的指标（很快就会更充分地解释），而（微分形式）是某个示性类的规范不变微分形式表示的积分。

当 $M_{d+1}=\varnothing$、即 $N_{d+2}$ 闭时，这就是 Atiyah--Singer 指标定理，它把 Dirac 算子的指标等同于某个示性类的积分。Atiyah--Patodi--Singer 指标定理是它对带边流形的推广，而 eta 不变量最初正是作为与边界相关的修正项被引入的。在本讲义中，我们反过来独立地引入 eta 不变量，把它当作质量无穷大的自由费米子理论配分函数的相位部分。

在作更一般的讨论之前，先给一个例子，其中 (12.10) 里的三项都能各自显式算出。正如 [AC79] 最早指出的，它在相当初等的量子力学设定中也有出人意料自然的出现。\footnote{它也曾被用作我们研究生院硕士课程入学考试 2010 年代的笔试题之一。}

我们知道 Landau 能级的物理：二维电子在恒定磁场 $B$ 中运动。这里用 Pauli 方程把电子自旋的效应显式地包含进来，并考虑一个无限的 $(x,y)$ 平面，其中 $B_z(x,y)$ 集中在有限区域 $r<R$ 内，使得
\begin{equation}
  \Phi:=\int_D B_z(x,y)\,dx\,dy
  \tag{12.11}
\end{equation}
有限。我们预期大约有 $\sim|\Phi|$ 个简并零模被束缚在区域 $r<R$ 中。下面来显式验证这一点。

% [原文 p.174]
电子的波函数这时有两个分量：
\begin{equation}
  \vec\psi(x,y)=\begin{pmatrix}\psi_\uparrow(x,y)\\ \psi_\downarrow(x,y)\end{pmatrix},
  \tag{12.12}
\end{equation}
而哈密顿量为
\begin{equation}
  H=\Bigl[-i\frac{\partial}{\partial x}+A_x(x,y)\Bigr]^2
  +\Bigl[-i\frac{\partial}{\partial y}+A_y(x,y)\Bigr]^2
  +B_z(x,y)\sigma_z,
  \tag{12.13}
\end{equation}
这里我们取了 $m=\hbar=e=1$ 的单位制，且 $B_z=\partial_xA_y-\partial_yA_x$。注意 $H$ 与 $\sigma_z$ 平凡地交换，所以自旋向上与自旋向下的电子各自独立运动。

话虽如此，考虑下面这个算子会更方便：
\begin{equation}
  D=\sigma_x\Bigl[-i\frac{\partial}{\partial x}+A_x(x,y)\Bigr]
  +\sigma_y\Bigl[-i\frac{\partial}{\partial y}+A_y(x,y)\Bigr],
  \tag{12.14}
\end{equation}
它与 $\sigma_z$ 反交换，并满足
\begin{equation}
  H=D^2.
  \tag{12.15}
\end{equation}
由于 $D$ 是 Hermite 的，$H=D^2$ 表明 $H$ 的本征值 $E$ 满足 $E\ge0$；而波函数 $\psi$ 是 $H$ 的属于 $E=0$ 的本征矢，当且仅当 $D\psi=0$。

现在设 $E\neq0$，取 $H$ 的本征矢 $\psi$ 使 $H\psi=E\psi$ 且 $\sigma_z\psi=+\psi$。令 $\psi':=D\psi\neq0$。容易验证 $\sigma_z\psi'=-\psi'$ 且 $H\psi'=E\psi'$。因此只要 $E\neq0$，$H$ 的 $\sigma_z=+1$ 与 $\sigma_z=-1$ 的本征矢就成对出现。

相反，$E=0$ 时没有这种配对，因为 $D\psi=0$。我们把 $E=0$ 的本征态再仔细研究一下。

现在引入函数 $\phi(x,y)$，使其满足
\begin{equation}
  A_x(x,y)=\frac{\partial\phi}{\partial y},\qquad
  A_y(x,y)=-\frac{\partial\phi}{\partial x}.
  \tag{12.16}
\end{equation}
于是 $B_z=-\Delta\phi$。关系 (12.11) 意味着，当 $r\gg0$ 时
\begin{equation}
  \phi(x,y)\sim\frac{\Phi}{2\pi}\log r .
  \tag{12.17}
\end{equation}
确实，
\begin{equation}
  \int B_z\,dx\,dy=\int_{S^1\text{ at large }R}A
  =\int_0^{2\pi}(\operatorname{div}\phi\cdot\vec n)\,d\theta
  =2\pi R\,\frac{\Phi}{2\pi}\,\frac{1}{R}=\Phi .
  \tag{12.18}
\end{equation}

为解 $D\psi=0$，我们设 $\sigma_z=+1$ 或 $\sigma_z=-1$。再通过
\begin{equation}
  \psi_\uparrow(x,y)=f_\uparrow(x,y)e^{\phi(x,y)},\qquad
  \psi_\downarrow(x,y)=f_\downarrow(x,y)e^{-\phi(x,y)}
  \tag{12.19}
\end{equation}
定义 $f_\uparrow(x,y)$ 与 $f_\downarrow(x,y)$。

% [原文 p.175]
方程 $D\psi=0$ 容易改写成
\begin{equation}
  \frac{\partial}{\partial\bar z}f_\uparrow=0,\qquad
  \frac{\partial}{\partial z}f_\downarrow=0,
  \tag{12.20}
\end{equation}
其中 $z=x+iy$、$\bar z=x-iy$。这意味着 $f_\uparrow$ 是全纯函数，$f_\downarrow$ 是反全纯函数。于是
\begin{equation}
  \psi_\uparrow=P_\uparrow(z)e^{\phi},\qquad
  \psi_\downarrow=P_\downarrow(\bar z)e^{-\phi},
  \tag{12.21}
\end{equation}
其中 $P_\uparrow(z)$ 与 $P_\downarrow(\bar z)$ 是多项式。

注意 $e^{\phi}\sim r^{\Phi/2\pi}$，$e^{-\phi}\sim r^{-\Phi/2\pi}$。为确定起见，取 $\Phi$ 为负。此时 $\psi_\downarrow$ 永远不可平方可积；而 $\psi_\uparrow$ 平方可积的条件是多项式
\begin{equation}
  P_\uparrow(z)=c_nz^n+c_{n-1}z^{n-1}+\cdots
  \tag{12.22}
\end{equation}
的最高非零项 $z^n$ 满足 $n\le(\Phi/2\pi)-1$。于是我们发现，$\psi_\uparrow$ 构成的解空间是 $\lfloor\Phi/2\pi\rfloor$ 维的。

总结一下：$\sigma_z=\pm1$ 的 $E=0$ 解的数 $N_\pm$ 由
\begin{equation}
  N_+=\Bigl\lfloor\frac{\Phi}{2\pi}\Bigr\rfloor,\qquad N_-=0
  \tag{12.23}
\end{equation}
给出。

我们把它写成
\begin{equation}
  \tilde\eta(a)=(N_+-N_-)-\int\frac{F}{2\pi},
  \tag{12.24}
\end{equation}
其中
\begin{equation}
  a=\int_{S^1\text{ at large }r=R}A
  \tag{12.25}
\end{equation}
而
\begin{equation}
  \tilde\eta(a)=\lfloor a\rfloor-a .
  \tag{12.26}
\end{equation}
我们的结果 (12.24) 应当与一般的 APS 指标定理 (12.10) 相比较：

\begin{itemize}
  \item $N_+-N_-$ 是体 Dirac 算子 (12.14) 的指标，那里的“手征性”就是 $\sigma_z$ 本征值；
  \item $\int\frac{F}{2\pi}$ 是 $U(1)$ 丛的第 1 Chern 类的微分形式表达式；
  \item $-\tilde\eta(a)$ 是在非常大的半径处求值的 $S^1$ 上的 eta 不变量，只是少了减去 $1/2$ 这一步。少了这一步，是因为 APS 定理应当在带柱状边界的流形上表述。这时需要作从 $(x,y)$ 到 $(r,\theta)$ 的坐标变换，而这一变换会交换费米子的周期与反周期边界条件，这一点几乎让每个学 2 维 CFT 与弦论的学生都困惑过。
\end{itemize}


\subsection{Atiyah--Patodi--Singer 指标定理：一般讨论}

% [原文 p.176]
\para{奇数维的 $M_{d+1}$ 与偶数维的 $N_{d+2}$。}
下面在一般情形下解释 APS 指标定理 (12.10) 中出现的各项。先从边界 $M_{d+1}$ 为奇数维的情形说起。这时考虑 $N_{d+2}$ 上的 Dirac 算子 $D_N$，它作用在复旋量上，而 $N_{d+2}$ 是偶数维的。这个旋量不仅可以耦合到时空曲率，还可以耦合到 $U(k)$ 规范场 $F$。

在偶数维，有正手征与负手征的旋量，由与 $D_N$ 反交换的手征算子 $\Gamma$ 的本征值区分。令 $n_\pm$ 为满足 $D\psi=0$ 且 $\Gamma\psi=\pm\psi$ 的零模个数，则 $D_N$ 的指标定义为
\begin{equation}
  \operatorname{index}D_N:=n_+-n_- .
  \tag{12.27}
\end{equation}
已知微分形式的贡献取下形：
\begin{equation}
  \int_N(\text{微分形式})=\int_N\wh{A}(N)\cdot\tr e^{F/(2\pi)},
  \tag{12.28}
\end{equation}
其中 $\wh{A}(N)$ 是由空间曲率构造的微分形式，称为 $\wh{A}$ 亏格；$F$ 是 $U(k)$ 曲率，我们把它取为取值在 $k\times k$ Hermite 矩阵中的微分形式。$\wh{A}(N)$ 的显式为
\begin{equation}
  \wh{A}(N)=1-\frac{1}{24}p_1+\frac{1}{5760}\bigl(-4p_2+7p_1^2\bigr)+\cdots,
  \tag{12.29}
\end{equation}
其中
\begin{equation}
  p_1=-\frac{1}{2}\tr\Bigl(\frac{R}{2\pi}\Bigr)^2,\qquad
  p_2=\frac{1}{2}\Bigl(\Bigl(\frac{1}{2}\tr\Bigl(\frac{R}{2\pi}\Bigr)^2\Bigr)^2
  -\frac{1}{2}\tr\Bigl(\frac{R}{2\pi}\Bigr)^4\Bigr),\ \cdots
  \tag{12.30}
\end{equation}
是 Pontryagin 类。

例如，当边界 $M_3$ 是三维的、体 $N_4$ 是四维的，且旋量耦合到 $U(1)$ 规范场 $F$ 时，有
\begin{equation}
  \eta(D_M)=\operatorname{index}D_N
  +\int_N\Bigl(\frac{1}{2}c_1^2-\frac{1}{24}p_1\Bigr),
  \tag{12.31}
\end{equation}
其中 $c_1=F/2\pi$ 就是我们前阵子讨论过的第 1 Chern 类。我们也可以考虑 $U(1)$ 规范场为零的几乎同样的情形。用上标 $0$ 记对应的 Dirac 算子，则
\begin{equation}
  \eta(D_M^0)=\operatorname{index}D_N^0+\int_N\Bigl(-\frac{1}{24}p_1\Bigr).
  \tag{12.32}
\end{equation}
于是
\begin{equation}
  \eta(D_M)-\eta(D_M^0)
  =\bigl(\operatorname{index}D_N-\operatorname{index}D_N^0\bigr)
  +\int_N\frac{1}{2}c_1^2 .
  \tag{12.33}
\end{equation}
这意味着 $\eta(D_M)-\eta(D_M^0)$ 本质上是一个（更好定义的）阿贝尔 Chern--Simons 项
\begin{equation}
  \frac{1}{2}\,\frac{A}{2\pi}\,\frac{F}{2\pi},
  \tag{12.34}
\end{equation}
% [原文 p.177]
它在模 $1$ 意义下良定义。回想 $c_1$ 是 $H^2(N;\Z)$ 中的自然元素，$c_1^2$ 自然落在 $H^4(N;\Z)$ 中。让除以 $1/2$ 这件事有意义，曾是 §8.2 的一大主题。这里我们从另一个角度看到了它。

由式 (12.32) 可以把 $\eta(D_M^0)$ 看作一个更好定义的引力 Chern--Simons 项，参看 §8.3.3 的讨论。对一般的 $SO(n)$ 丛，$p_1\in H^4(N;\Z)$ 是自然整值定义的上同调类，其微分形式代表如上。由式 (12.32) 可见
\begin{equation}
  \int_N\frac{p_1}{24}=\operatorname{index}D_N\in\Z .
  \tag{12.35}
\end{equation}
其实，当 $d+2=4$ 时，旋量丛有四元数结构，即存在反线性映射 $\psi\mapsto j\psi$，它与 $D_N$ 及 $\Gamma$ 都交换。（这是 Kramers 简并的欧氏版本。）因此 $D_N$ 与 $\Gamma$ 的所有本征空间都是二重简并的，$n_\pm$ 都是偶数，从而 $\operatorname{index}D_N$ 是偶数。于是实际上有
\begin{equation}
  \int_N\frac{p_1}{48}\in\Z .
  \tag{12.36}
\end{equation}

记 spin 联络为 $\omega$，并记 $R=d\omega+\omega\wedge\omega$，则
\begin{equation}
  d\tr\Bigl(\omega d\omega+\frac{2}{3}\omega^3\Bigr)=\tr R^2 .
  \tag{12.37}
\end{equation}
因此 $\eta(D_M)$ 可以看作
\begin{equation}
  \frac{1}{48}\int_M\frac{1}{2\cdot(2\pi)^2}
  \tr\Bigl(\omega d\omega+\frac{2}{3}\omega^3\Bigr)
  \tag{12.38}
\end{equation}
的更好版本，后者在模 $1$ 意义下良定义。多出来的因子 $1/48$ 是最不平凡的部分。

回到 eta 不变量的定义：它是 $2+1$ 维中单个质量极大的费米子的有效作用量的相位部分。对这一有效作用量套用 Kubo 公式，就得到具有最小量子化系数的热量子霍尔效应。

\para{偶数维的 $M_{d+1}$ 与奇数维的 $N_{d+2}$。}
当边界 $M_{d+1}$ 是偶数维时，微分形式的贡献直接为零：由曲率 2 形式构造的微分形式总是偶次，它们在奇数维的 $N_{d+2}$ 上的积分为零。手征算子需要有几分小心地定义，但最终指标的定义与前面一样，于是得到更简单的公式
\begin{equation}
  \eta(D_M)=\operatorname{index}D_N .
  \tag{12.39}
\end{equation}
这意味着，偶数维流形若是边界、即形如 $M_{d+1}=\partial N_{d+2}$，其上的 eta 不变量模 $1$ 为零。换言之，$e^{2\pi i\eta(D_M)}=1$。注意这只对是边界的 $M$ 成立。这一点很快会起重要作用。

\subsection{另一个例子：$T^2$、$S^2$ 与 $S^2/\Z_2=\RP^2$}

% [原文 p.178]
\para{$T^2$ 上的复大质量费米子}
作为另一组例子，我们来看二维的 eta 不变量。最简单的情形是 $T^2$。考虑带 $U(1)$ 荷的复旋量，这就给出与 §12.1.3 中跟随 Aharonov--Casher 考虑过的同一个算子 $D$，只不过把无限的 $(x,y)$ 平面换成了有限的周期环面。但要注意，那里我们把二维部分当作体、把无穷远处的 $S^1$ 当作边界；这里考虑的是定义在二维部分上的 $D$ 的 eta 不变量。

不过那里的分析可以照搬。$H=D^2$ 的任何非零本征值都成对出现，两个本征值的 $\Gamma=\sigma_z$ 相反。在这个基下 $D$ 形如 $\begin{pmatrix}0&x\\ x&0\end{pmatrix}$，从而 $D$ 的本征值为 $\pm x$。只要 $x\neq0$，这一对本征值对 eta 不变量的贡献就是 $(+1)e^{-s|x|}+(-1)e^{-s|x|}$，在取 $s\to0$ 极限之前就已经相消。所以只有零本征值有贡献，于是
\begin{equation}
  \eta(D_{T^2})=\frac{n_++n_-}{2}\mod 1,
  \tag{12.40}
\end{equation}
其中 $n_\pm$ 是 $D$ 的满足 $\Gamma=\sigma_z=\pm1$ 的零模个数。

利用指标定理，我们知道
\begin{equation}
  \operatorname{index}D_{T^2}=n_+-n_-
  =\int_{T^2}\frac{F}{2\pi}=:n .
  \tag{12.41}
\end{equation}
（当 $F$ 均匀时，$(n_+,n_-)=(|n|,0)$ 或 $(0,|n|)$，取决于 $n$ 的符号。这就是标准的 Landau 简并。）由于 $(n_++n_-)/2\equiv(n_+-n_-)/2\bmod 1$，我们得到
\begin{equation}
  \eta(D_{T^2})=\frac{n}{2}\mod 1,
  \tag{12.42}
\end{equation}
等价地
\begin{equation}
  e^{2\pi i\eta(D_{T^2})}=(-1)^n .
  \tag{12.43}
\end{equation}

\para{$T^2$ 上的实大质量费米子}
当 $n=\int_{T^2}F/2\pi=0$ 时，$e^{2\pi i\eta(D_{T^2})}=+1$。但注意这是欧氏号差下复的大质量费米子的结果。

若从二维的实费米子出发并 Wick 旋转到欧氏号差，我们仍需要把自由度减半，尽管二维欧氏号差下的旋量是复的、不能考虑实值旋量场。结果欧氏 Dirac 谱的本征值自动二重简并，\footnote{取 $D=(i\partial_x)\sigma_x+(i\partial_y)\sigma_y$，$T=K\sigma_y$，其中 $K$ 是两个分量旋量上的复共轭。这个 $T$ 与 $(x,y)$ 平面的 $\theta$ 旋转交换，因此可以一致地作用在弯曲流形上的旋量上。我们有 $T^2=-1$ 且 $DT=TD$。这个代数迫使 $D$ 的每个本征值二重简并；所以这就是标准的 Kramers 简并。还要注意 $R_x:\psi(x,y)\mapsto\sigma_z\sigma_x\psi(-x,y)$ 与 $R_y:\psi(x,y)\mapsto\sigma_z\sigma_y\psi(x,-y)$ 与 $D$ 交换，并在费米子波函数上定义一个 pin$^-$ 作用。它们也与 $T$ 交换，所以在不可定向的 pin$^-$ 流形上本征态同样是二重简并的。归根结底这些都来自 $\Spin(2)\subset\Pin^-(2)\subset\Spin(3)$ 这一事实，以及 $\Spin(3)$ 的旋量表示具有自然的四元数结构。}我们可以一致地取本征值的一半。

% [原文 p.179]
例如在 $n=0$ 的 $T^2$ 上，我们就只有 $T^2$ 上的周期旋量。显然 $D$ 有两个零模。取它们的一半，就只剩一个零本征值，于是
\begin{equation}
  \eta\bigl(D^{\text{real}}_{T^2}\bigr)=\frac{1}{2},
  \tag{12.44}
\end{equation}
并且
\begin{equation}
  e^{2\pi i\eta(D^{\text{real}}_{T^2})}=-1 .
  \tag{12.45}
\end{equation}

\para{$S^2$ 上的实大质量费米子}
接着考虑 $S^2$ 上的 Dirac 算子。我们都知道 $S^2$ 上标量值函数如何分解为 $so(3)$ 的表示：有 $V_\ell$，$\ell=0,1,2,\dots$，$\dim V_\ell=2\ell+1$。旋量值函数可以通过张量 $V_{1/2}$ 得到。结果对 $\ell=1/2,3/2,\dots$ 各有两份 $V_\ell$。$D$ 在一份 $V_\ell$ 上的本征值为 $2\ell+1$，在另一份上的本征值为 $-(2\ell+1)$。注意如前所述，每个本征值都出现偶数次。无论如何，eta 不变量自动为零。

\para{$\RP^2$ 上的实大质量费米子}
再考虑 $S^2/\Z_2=\RP^2$。它不可定向，但有 pin$^-$ 结构。为看出这一点，把 $\RP^2$ 看作北半球加上赤道上的认同。赤道附近的 spin 结构应当是 NS（反周期）的，因为赤道可以缩成绕北极的小圆。而绕赤道一周相当于做两次对径认同。由于对径认同会翻转定向，这意味着定向翻转作用在旋量上平方为 $-1$。所以这是 pin$^-$。

更细致的分析表明，在 $\RP^2$ 上，(s)pinor 波函数对每个 $V_\ell$、$\ell=1/2,3/2,\dots$ 都分解为单份，而 Dirac 本征值带有交错符号 $(-1)^{\ell-1/2}(2\ell+1)$。于是实费米子的 eta 不变量为
\begin{equation}
  \eta\bigl(D^{\text{real}}_{\RP^2}\bigr)
  =\lim_{s\to0}\frac{1}{2}\sum_{\ell=\frac12,\frac32,\dots}
  (-1)^{\ell-1/2}\frac{2\ell+1}{2}e^{-s(2\ell+1)}=\frac{1}{8},
  \tag{12.46}
\end{equation}
其中分母 $1/2$ 来自：为考虑实费米子的配分函数，需要把 Dirac 本征值取一半。

另一种用 APS 指标定理的计算如下。考虑 $T^3$，坐标为 $(x,y,z)$，其中 $x\sim x+1$、$y\sim y+1$、$z\sim z+1$。取商 $(x,y,z)\sim(-x,-y,-z)$。在 $x,y,z\in\{0,1/2\}$ 处有 $2^3$ 个不动点。每个不动点的邻域看起来像 $\R^3/\Z_2$。于是在每个不动点周围挖去一个以 $\RP^2$ 为边界的 $B^3/\Z_2$，就得到光滑流形 $M_3$，其边界是八个 $\RP^2$。pinor 丛在 $M_3$ 上是二维的，Dirac 零模就是常数，共有两个零模。由于我们关心的是实费米子，需要从这两个里取一个。因此
\begin{equation}
  8\eta\bigl(D^{\text{real}}_{\RP^2}\bigr)=1,
  \tag{12.47}
\end{equation}
从而
\begin{equation}
  \eta\bigl(D^{\text{real}}_{\RP^2}\bigr)=1/8 .
  \tag{12.48}
\end{equation}

\takeaway{本节把 eta 不变量定义为质量无穷大的自由费米子配分函数的相位部分，即谱中本征值符号的正规化和的一半。它可写成 $e^{2\pi i\eta(D)}=Z$。在 $S^1$ 与 Landau 能级的例子中，直接求和与 APS 指标定理给出同一结果。}


\section{可逆相的一般低能有效作用量}

% [原文 p.180]
到此为止，我们弄清了 eta 不变量的性质，也就是质量极大的自由费米子的低能有效作用量的相位部分。现在把它外推，对一般可逆相的相位部分作出公设。这里的可逆相指这样的有能隙的相：在闭空间流形上，其基态希尔伯特空间总是一维的。

给定一个 $d+1$ 维的理论 $T$，把它放到 Wick 旋转后的欧氏流形 $M_{d+1}$ 上。该流形带有度规；若理论 $T$ 有 $G$ 对称性，还带有背景 $G$ 规范场。例如，若 $T$ 有 $U(1)$ 对称性（对应电磁相互作用），就给 $M_{d+1}$ 配上一个 $U(1)$ 规范场，即我们用来探测系统 $T$ 的 Maxwell 场。为记号简洁，我们把这些附加数据隐含不写，只把整体记作 $M_{d+1}$。

于是就有配分函数，即有效作用量的指数，其形式为
\begin{equation}
  Z_T(M_{d+1})=\exp\bigl(-S_T(M_{d+1})\bigr).
  \tag{12.49}
\end{equation}
我们关注 $Z_T$ 的相位部分，即 $S_T$ 的虚部，并把它写成
\begin{equation}
  S_T(M_{d+1})=\text{实部}+2\pi i\,I_T(M_{d+1}).
  \tag{12.50}
\end{equation}
这里 $I_T(M_{d+1})$ 应当在模 $1$ 意义下良定义。

我们假设以下几条性质：

\begin{enumerate}
  \item 对不交并，应当有
  \begin{equation}
    I_T(M_{d+1}\sqcup M'_{d+1})
    =I_T(M_{d+1})+I_T(M'_{d+1}).
    \tag{12.51}
  \end{equation}
  \item 在定向反转下，有
  \begin{equation}
    I_T(\ol{M_{d+1}})=-I_T(M_{d+1}).
    \tag{12.52}
  \end{equation}
  \item 存在一个由度规曲率 $R$ 与 $G$ 曲率 $F$ 构造的、次数为 $d+2$ 的示性类 $J_T(R,F)$，使得当 $M_{d+1}$ 是边界、$M_{d+1}=\partial N_{d+2}$ 时，有
  \begin{equation}
    I_T(M_{d+1})=\int_{N_{d+2}}J_T(R,F).
    \tag{12.53}
  \end{equation}
\end{enumerate}

对每条性质作些评论：

\begin{enumerate}
  \item 第一条性质直截了当。
  \item 第二条性质是 Wick 旋转后作用量的一般性质，可以称为“反射实性”，尽管这个术语用得并不广。（常用的是“反射正性”，它在这一性质之上再加一条正性约束。）

  一个简短的说明如下。（有效）作用量在 Wick 旋转之前是实的。现在把 $t$ 换成 $\tau=it$。含偶数个 $t$ 导数的项保持为实，含奇数个 $t$ 导数的项则变为虚。唯一具有相对论不变性的这类项正比于 $\epsilon_{tx_1x_2\cdots x_d}$，而它在定向反转下是奇的。
% [原文 p.181]
  \item 第三条性质最不平凡。它的一个推论如下。取定一个 $M_{d+1}$，固定度规 $g_{\mu\nu}$，并考虑背景电磁势 $A_\mu$。作一个小变分 $A'_\mu=A_\mu+\delta A_\mu$。我们想知道低能有效作用量虚部的变分，即 $I_T(M_{d+1},A'_\nu)-I_T(M_{d+1},A_\nu)$。为此取 $N_{d+1}:=M_{d+1}\times[0,1]$，其边界为
  \[
    \partial N_{d+1}=M_{d+1}|_{\text{at }1}
    \sqcup\ol{M_{d+1}}|_{\text{at }0}.
  \]
  可以在 $N_{d+1}$ 上配置度规与 Maxwell 场，使其在 $0$ 处限制为 $A_\nu$，在 $1$ 处限制为 $A'_\nu$。于是
  \begin{equation}
    I_T(M_{d+1},A'_\nu)-I_T(M_{d+1},A_\nu)
    =\int_{M_{d+1}\times[0,1]}J_T(R,F).
    \tag{12.54}
  \end{equation}

  现在取 $d=2$、$J_T(R,F)=aF^2$ 为例。由于 $aF^2=ad(AF)$，可以分部积分，得到
  \begin{equation}
    I_T(M_{d+1},A'_\nu)-I_T(M_{d+1},A_\nu)
    =a\int_{M_{d+1}}A'F'-a\int_{M_{d+1}}AF .
    \tag{12.55}
  \end{equation}

  所以第三条性质意味着：$I_T(M_{d+1})$ 的小变分等于与示性类 $J_T(R,F)$ 相伴的 Chern--Simons 项的小变分，等等。

  但第三条性质还要更强，因为 $N_{d+2}$ 并不限于 $M_{d+1}\times[0,1]$ 这种形式。
\end{enumerate}

eta 不变量，即质量极大的费米子的低能有效作用量的相位部分，确实满足这三条性质。我们假设任意相互作用系统的低能有效作用量的相位部分也满足这三条性质，并采取这样的观点：分类这种相互作用可逆相，就是分类满足这三条性质的 $I(M_{d+1})$。

\takeaway{可逆相的低能有效作用量的相位部分写成 $2\pi iI_T$，并公设三条：不交并相加、定向反转取负、在边界上等于示性类 $J_T$ 的积分。第三条等价于说 $J_T$ 给出配边不变量，且这一要求比局部 Chern--Simons 项更强。}

\section{配边群}

为此，配边群这一代数拓扑中的另一重要概念将起重要作用。下面就来介绍它。

\subsection{定义}

最朴素的版本如下。选定一个想加在所要考虑的流形上的结构 $\mathcal S$。（$\mathcal S$ 可以只是定向，也可以是 spin 结构、pin$^\pm$，或者只是光滑结构，等等。）给定一个带若干边界分支的流形，我们区分它的入边界与出边界，记作
\begin{equation}
  \partial N_{d+1}=M^{out}_d\sqcup\ol{M^{in}_d}.
  \tag{12.56}
\end{equation}
我们称 $N_{d+1}$ 为从 $M^{in}_d$ 到 $M^{out}_d$ 的一个配边，并记作
\begin{equation}
  N_{d+1}:M^{in}_d\to M^{out}_d .
  \tag{12.57}
\end{equation}
这时称 $M^{in}_d$ 与 $M^{out}_d$ 配边。若要强调 $\mathcal S$，就说 $\mathcal S$-配边（$\mathcal S$-bordism）、$\mathcal S$-配边的（$\mathcal S$-bordant），等等。

% [原文 p.182]
我们允许 $M^{in,out}_d$ 为 $\varnothing$。这时称 $N_{d+1}$ 为零配边。例如，若 $M^{out}_d=\varnothing$，就称 $M^{in}_d$ 是零配边的。

配边是一种等价关系。按配边分出的等价类的集合就是配边群
\begin{equation}
  \Omega^{\mathcal S}_d:=
  \frac{\bigl\{\text{$d$ 维带 $\mathcal S$ 结构的流形}\bigr\}}
       {\text{$\mathcal S$-配边}} ,
  \tag{12.58}
\end{equation}
其阿贝尔群运算由不交并给出：
\begin{equation}
  [M_d]+[M'_d]:=[M_d\sqcup M'_d].
  \tag{12.59}
\end{equation}
零元就是 $[\varnothing]$，而加法的逆是
\begin{equation}
  -[M_d]:=[\ol{M_d}],
  \tag{12.60}
\end{equation}
因为很容易造出一个配边
\begin{equation}
  M_d\sqcup\ol{M_d}\to\varnothing
  \tag{12.61}
\end{equation}
只需令 $N_{d+1}:=M_d\times[0,1]$ 即可。

我们还会考虑以下变体。作为结构，可以取原来的结构 $\mathcal S$ 再加一个到目标空间 $X$ 的映射。这时考虑数对 $(M_d,f)$，其中 $f:M_d\to X$。对 $(M_d,f)$ 与 $(M'_d,f')$ 之间的配边，可以要求体上的映射 $f:N_{d+1}\to X$ 在边界上限制为 $f:M_d\to X$ 与 $f':M'_d\to X$，等等，由此定义。取等价类，得到的群是
\begin{equation}
  \Omega^{\mathcal S}_d(X):=
  \frac{\bigl\{\text{$d$ 维带 $\mathcal S$ 结构的流形与到 $X$ 的映射}\bigr\}}
       {\text{配边}} .
  \tag{12.62}
\end{equation}
特别地我们看到
\begin{equation}
  \Omega^{\mathcal S}_d(\pt)=\Omega^{\mathcal S}_d\quad(\text{即上面定义的}).
  \tag{12.63}
\end{equation}
另外还有
\begin{equation}
  \Omega^{\mathcal S}_d(BG)=
  \frac{\bigl\{\text{$d$ 维带 $\mathcal S$ 结构的流形与一个 $G$ 丛}\bigr\}}
       {\text{配边}},
  \tag{12.64}
\end{equation}
因为映射 $f:M_d\to BG$ 分类 $M_d$ 上的一个 $G$ 丛。

关于术语的一点评论。在较早的数学文献中，更常见的说法是 cobordism、cobordant 等等；而在较新的文献中，说 bordism、bordant 等等正变得越来越常见。原因似乎是人们后来意识到，群 $\Omega^{\mathcal S}_d(X)$ 是一般化同调的例子，而不是一般化上同调。这大概就导致把词首的前缀 co 去掉了。

\subsection{例子与一般性质}

对足够好的空间 $X$，$\Omega^{\mathcal S}_d(X)$ 是有限生成的，因此形如
\begin{equation}
  \Omega^{\mathcal S}_d(X)=\Z^k\oplus\bigoplus_i\Z_{n_i}.
  \tag{12.65}
\end{equation}
$\Z^k$ 这一部分称为自由部分或非挠部分，其余部分称为挠部分。我们常遇到的配边群汇总在表 10 中。

下面就从这张表来了解这些配边群的一些性质。

% [原文 p.183]
\begin{center}
  \small
  \setlength{\tabcolsep}{4pt}
  \begin{tabular}{c||cc|ccc|ccc}
    \toprule
    $d$
    & $\Omega_d^{\mathrm{unoriented}}$ & $\Omega_d^{\mathrm{oriented}}$
    & $\Omega_d^{\mathrm{pin}^-}$ & $\Omega_d^{\mathrm{spin}}$ & $\Omega_d^{\mathrm{pin}^+}$
    & $\Omega_d^{\mathrm{spin}}(B\Z_2)$ & $\Omega_d^{\mathrm{spin}}(BU(1))$
    & $\Omega_d^{\mathrm{spin}\text{-}c}$ \\
    \midrule
    0 & $\Z_2$ & $\Z$ & $\Z_2$ & $\Z$ & $\Z_2$ & $\Z$ & $\Z$ & $\Z$ \\
    1 & $0$ & $0$ & $\Z_2$ & $\Z_2$ & $0$ & $\Z_2^2$ & $\Z_2$ & $0$ \\
    2 & $\Z_2$ & $0$ & $\Z_8$ & $\Z_2$ & $\Z_2$ & $\Z_2^2$ & $\Z\oplus\Z_2$ & $\Z$ \\
    3 & $0$ & $0$ & $0$ & $0$ & $\Z_2$ & $\Z_8$ & $0$ & $0$ \\
    4 & $\Z_2^2$ & $\Z$ & $0$ & $\Z$ & $\Z_{16}$ & $\Z$ & $\Z^2$ & $\Z^2$ \\
    5 & $\Z_2$ & $\Z_2$ & $0$ & $0$ & $0$ & $0$ & $0$ & $0$ \\
    6 & $\Z_2^3$ & $0$ & $\Z_{16}$ & $0$ & $0$ & $0$ & $\Z^2\oplus\Z_2$ & $\Z^2$ \\
    \bottomrule
  \end{tabular}

  \smallskip
  {\small 表 10：常遇到的配边群。（还没有仔细核对过，所以请预料到其中可能有笔误……）\par}
\end{center}

\para{约化配边群}
一个有用的事实是，我们总有直和分解
\begin{equation}
  \Omega^{\mathcal S}_d(X)=\Omega^{\mathcal S}_d(\pt)\oplus
  \wt\Omega^{\mathcal S}_d(X),
  \tag{12.66}
\end{equation}
其中 $\wt\Omega^{\mathcal S}_d(X)$ 称为约化配边群。在上面的表中，$\Omega^{\mathrm{spin}}(B\Z_2)$ 与 $\Omega^{\mathrm{spin}}(BU(1))$ 两列就能看到这一分解。

这一分解可以通过构造两个同态
\begin{equation}
  \Omega^{\mathcal S}_d(X)\to\Omega^{\mathcal S}_d(\pt),\qquad
  \Omega^{\mathcal S}_d(\pt)\to\Omega^{\mathcal S}_d(X)
  \tag{12.67}
\end{equation}
来证明，它们是互逆的。做法很简单：把 $[M,f]$ 送到 $[M]$，把 $[M]$ 送到 $[M,c]$，其中 $c:M\to X$ 是到某点 $p\in X$ 的常值映射。

\para{实值示性类}
实值示性类按以下意义给出配边不变量。例如取
\begin{equation}
  p_1=-\frac{1}{2}\tr\Bigl(\frac{R}{2\pi}\Bigr)^2 .
  \tag{12.68}
\end{equation}
给定配边 $N_5:M_4\to M'_4$，$p_1$ 可以定义在 $N_5$ 上，并自然地限制到 $M_4$ 与 $M'_4$。因此
\begin{equation}
  \int_{M'_4}p_1-\int_{M_4}p_1=\int_{N_5}dp_1=0,
  \tag{12.69}
\end{equation}
也就是说，$p_1$ 在两个配边流形上的积分相同。这就给出一个同态
\begin{equation}
  p_1:\Omega^{\mathcal S}_4\to\R .
  \tag{12.70}
\end{equation}
显然，它只能探测 $\Omega^{\mathcal S}_4$ 的自由部分 $\Z^k$。

% [原文 p.184]
反过来，已知 $\Omega^{\mathcal S}(X)$ 的所有自由部分都能用实值示性类探测。这就很容易数出 $\Z$ 直和项在 $\Omega^{\mathcal S}_d(X)$ 中出现多少份。例如，$\Omega^{\mathrm{spin}}_4(BU(1))\simeq\Z^2$ 可以用 $p_1$ 与
\begin{equation}
  c_1^2=\Bigl(\frac{F}{2\pi}\Bigr)^2
  \tag{12.71}
\end{equation}
来探测；而 $\Omega^{\mathrm{spin}}_6(BU(1))$ 的 $\Z^2$ 部分由 $p_1c_1$ 与 $c_1^3$ 探测。

确定这些配边不变量的对偶整基则是更精细的问题。例如 $\Omega^{\mathrm{oriented}}_4(\pt)$ 与 $\Omega^{\mathrm{spin}}_4(\pt)$ 都是 $\Z$，都能由 $\int_{N_4}p_1$ 探测。然而 $\Omega^{\mathrm{oriented}}_4(\pt)\simeq\Z$ 由 $\CP^2$ 生成，且
\begin{equation}
  \int_{\CP^2}\frac{p_1}{3}=1 ;
  \tag{12.72}
\end{equation}
而 $\Omega^{\mathrm{spin}}_4(\pt)\simeq\Z$ 由 $K3$ 生成，且
\begin{equation}
  \int_{K3}\frac{p_1}{48}=-1 .
  \tag{12.73}
\end{equation}
因此 $\Omega^{\mathrm{oriented},\mathrm{spin}}_4(\pt)\simeq\Z$ 的对偶基分别由 $p_1/3$ 与 $p_1/48$ 给出。换言之，由忘掉 spin 结构得到的同态
\begin{equation}
  \Z\simeq\Omega^{\mathrm{spin}}_4(\pt)\to
  \Omega^{\mathrm{oriented}}_4(\pt)\simeq\Z
  \tag{12.74}
\end{equation}
是乘以 $-16$。

类似地，$\Omega^{\mathrm{oriented}}_4(BU(1))$ 与 $\Omega^{\mathrm{spin}}_4(BU(1))$ 都是 $\Z$，都由 $c_1^2$ 探测，但对偶基分别是 $c_1^2$ 与 $(1/2)c_1^2$，而同态
\begin{equation}
  \Z\simeq\wt\Omega^{\mathrm{spin}}_4(BU(1))\to
  \wt\Omega^{\mathrm{oriented}}_4(BU(1))\simeq\Z
  \tag{12.75}
\end{equation}
是乘以 $2$。

\para{挠示性类}
带微分形式代表的实值示性类，在探测挠配边类时不太有用。这时可以改用 $\Z_p$ 值示性类，例如 Stiefel--Whitney 类。例如，给定流形 $M$，其切丛是定向向量丛，所以 Stiefel--Whitney 类 $w_i\in H^i(M;\Z_2)$ 自然有定义。它们给出配边不变量。例如，若有一个配边 $N_6:M_5\to M'_5$，则用 $\Z_2$ 值上链的 Stokes 定理有
\begin{equation}
  \int_{M'_5}w_2w_3-\int_{M_5}w_2w_3
  =\int_{N_6}\delta(w_2w_3)=0,
  \tag{12.76}
\end{equation}
这意味着有一个良定义的同态
\begin{equation}
  w_2w_3:\Omega^{\mathrm{oriented}}_5(\pt)\to\Z_2,
  \tag{12.77}
\end{equation}
% [原文 p.185]
而它确实已知能探测 $\Omega^{\mathrm{oriented}}_5(\pt)\simeq\Z_2$。（这个群的生成元是 $SU(3)/SO(3)$，即我们在上一节见过的对称 $3\times3$ 特殊酉矩阵的集合。）类似地，我们有良定义的同态
\begin{equation}
  w_2:\Omega^{\mathrm{unoriented}}_2(\pt)\to\Z_2,
  \tag{12.78}
\end{equation}
其生成元是 $\RP^2$。

\para{更精致的拓扑不变量}
示性类（即分类空间上 $\Z$ 值或 $\Z_p$ 值上同调类的拉回）不足以区分所有配边类。例如 $\Omega^{\mathrm{spin}}_1(\pt)\simeq\Z_2$ 由带周期 spin 结构的 $S^1$ 生成，但它并不由任何东西的积分给出：次数为 $1$ 的示性类只有 $w_1$，而 $w_1$ 探测的是切丛的不可定向性，$S^1$ 却是可定向的。

这类情形下 eta 不变量就有了用武之地。例如 $\eta(S^1)$ 在 spin 结构为周期（R）或反周期（NS）时分别是 $0$ 或 $1/2$ 模 $1$，从而探测 $\Omega^{\mathrm{spin}}_1(\pt)\simeq\Z_2$。类似地，$\eta(M_2)=0$ 或 $1/2$ 模 $1$ 探测 $\Omega^{\mathrm{spin}}_2(\pt)\simeq\Z_2$；后一个同态
\begin{equation}
  \eta:\Omega^{\mathrm{spin}}_2(\pt)\overset{\sim}{\longrightarrow}\Z_2
  \tag{12.79}
\end{equation}
也称为 Arf 不变量。它也正是 Kitaev 链非平凡相的配分函数的（相位或虚部）。

我们还提到过，$\Omega^{\mathrm{pin}^-}_2(\pt)\simeq\Z_8$ 由 eta 不变量探测：
\begin{equation}
  \eta:\Omega^{\mathrm{pin}^-}_2(\pt)\overset{\sim}{\longrightarrow}\Z_8,
  \tag{12.80}
\end{equation}
其中 $\eta(\RP^2)=1/8$ 模 $1$。

一般来说，并非所有配边群都能由实值示性类、$\Z_p$ 值示性类或 eta 不变量探测。但在实践中，在我们通常遇到的时空维数下，所有配边不变量都是这几类之一。

\para{如何确定配边群？}
确定配边群是一门手艺，并没有具体的算法能做到。例如 $\Omega^{\mathrm{spin}}_3(B\Z_2)$ 该如何求？

上面提到，非挠部分可由具有微分形式代表的实值示性类探测。这里没有这种东西，所以 $\Omega^{\mathrm{spin}}_3(B\Z_2)$ 是挠的。我们可以显式地计算 $S^3/\Z_2$ 上、带绕 $\pi_1(S^3/\Z_2)=\Z_2$ 非平凡 $\Z_2$ 丛的 $\Z_2$-奇费米子的 eta 不变量，结果等于 $1/8$。因此 $[S^3/\Z_2,f]$ 是 $\Omega^{\mathrm{spin}}_3(B\Z_2)$ 中的非平凡元素。现在，$n[S^3/\Z_2,f]=0$ 意味着 $n/8=0$ 模 $1$。把这样的最小正整数 $n$ 记作 $8k$。那么 $[S^3/\Z_2,f]$ 生成子群 $\Z_{8k}\subset\Omega^{\mathrm{spin}}_3(B\Z_2)$。又因为 $T^4/\Z_2$ 给出零配边 $2^4[S^3/\Z_2,f]=0$，所以 $k$ 只能是 $1$ 或 $2$。

还有一个方法叫 Atiyah--Hirzebruch 谱序列。由它可以相当容易地证明 $\#\Omega^{\mathrm{spin}}_3(B\Z_2)\le8$。合起来我们就知道 $\Omega^{\mathrm{spin}}_3(B\Z_2)\simeq\Z_8$。这类半凑合的分析往往管用，但需要时间、练习与经验。

% [原文 p.186]
其实 $\Omega^{\mathrm{spin}}_3(B\Z_2)$ 足够典型，另一种称为 Adams 谱序列的方法可以直接用来证明该配边群是 $\Z_8$。但要用 Adams 谱序列，得在代数拓扑上投入可观的精力。\footnote{Atiyah--Hirzebruch 谱序列要容易得多，背后还有一些物理上的理解。至于 Adams 谱序列的物理含义，据作者所知没有人知道。}

无论如何，当你需要知道某些配边群时，首先应当去查教科书或互联网，看是否已经有人算过。然后，应当去问你的朋友，或者干脆问我，看它能不能算、以及是否已知。其实大约一个月前，就有一个三位弦论学家组成的合作组这样联系过我，而我算出了 $\Omega^{\mathrm{spin}}_7(BG)\simeq0$，其中 $G=\prod_iG_i$，每个 $G_i$ 要么是 $U(1)$，要么是单连通紧致非阿贝尔单 Lie 群。

\takeaway{配边群按配边关系把带结构 $\mathcal S$ 的流形分类，群运算是不交并；带映射的版本记 $\Omega^{\mathcal S}_d(X)$。实值示性类只探测自由部分，挠部分要靠 $\Z_p$ 值示性类与 eta 不变量，一般情形没有算法，只能靠谱序列加经验。}

\section{配边与一般相互作用相}

\para{非挠部分}
回想一下：对于结构 $\mathcal S$ 与对称性 $G$，一个 $d+1$ 维的相互作用可逆相有一个次数为 $d+2$ 的实值示性类 $J_T(R,F)$，使得当 $M_{d+1}=\partial N_{d+2}$ 时
\begin{equation}
  I_T(M_{d+1})=\int_{N_{d+2}}J_T(R,F).
  \tag{12.81}
\end{equation}
注意这样的 $N_{d+2}$ 可能有多个。取另一个选择 $N'_{d+2}$。$I_T(M_{d+1})$ 必须在模 $1$ 意义下良定义，因此
\begin{equation}
  \int_{N'_{d+2}}J_T(R,F)-\int_{N_{d+2}}J_T(R,F)
  =\int_{N'_{d+2}\cup\ol{N_{d+2}}}J_T(R,F)\in\Z,
  \tag{12.82}
\end{equation}
其中 $N'_{d+2}\cup\ol{N_{d+2}}$ 是把 $N'_{d+2}$ 与（反向定向后的）$N_{d+2}$ 沿公共边界 $M_{d+1}$ 粘合得到的闭流形。取 $N'_{d+2}$ 为 $N_{d+2}\sqcup W_{d+2}$，其中 $W_{d+2}$ 是任意的闭流形（带结构 $\mathcal S$ 与 $G$ 丛），就发现
\begin{equation}
  \int_{W_{d+2}}J_T(R,F)\in\Z
  \tag{12.83}
\end{equation}
对所有 $W_{d+2}$ 成立，这对我们的目的是充分且必要的。这意味着 $J_T$ 是整值配边不变量：
\begin{equation}
  J_T:\Omega^{\mathcal S}_{d+2}(BG)\to\Z,
  \tag{12.84}
\end{equation}
或者等价地
\begin{equation}
  J_T\in\Hom\bigl(\Omega^{\mathcal S}_{d+2}(BG),\Z\bigr).
  \tag{12.85}
\end{equation}
当 $\Omega^{\mathcal S}_{d+2}(BG)\simeq\Z^k\oplus\bigoplus_i\Z_{n_i}$ 时，
\begin{equation}
  \Hom\bigl(\Omega^{\mathcal S}_{d+2}(BG),\Z\bigr)\simeq\Z^k .
  \tag{12.86}
\end{equation}
（由于这个原因，它们常被称为可逆相的自由部分；但这里的“自由”是阿贝尔群意义上的自由，而这里讨论的是一般非自由的、有相互作用的相。）

% [原文 p.187]
它们描述（一般化的）量子化的热霍尔电导与电霍尔电导。例如我们已提到，$\Omega^{\mathrm{oriented}}_4(BU(1))$ 与 $\Omega^{\mathrm{spin}}_4(BU(1))$ 都 $\simeq\Z^2$，但整基不同，这一点前面讨论过。前者即定向（=玻色）情形为
\begin{equation}
  J_T=\nu'\frac{p_1}{3}+n'c_1^2,
  \tag{12.87}
\end{equation}
后者即 spin（=费米）情形为
\begin{equation}
  J_T=\nu\frac{p_1}{48}+\frac{n}{2}c_1^2,
  \tag{12.88}
\end{equation}
其中 $\nu'$、$n'$、$\nu$、$n$ 都是整数。如前面所讨论的，它们对应量子化的热霍尔电导与电霍尔电导。

\para{挠部分}
对 $J_T$ 为零的可逆相，当存在配边 $N_{d+2}:M_{d+1}\to M'_{d+1}$ 时，有
\begin{equation}
  I_T(M'_{d+1})-I_T(M_{d+1})
  =\int_{N_{d+2}}J_T(R,F)=0,
  \tag{12.89}
\end{equation}
这意味着 $I_T$ 本身就是配边不变量
\begin{equation}
  I_T:\Omega^{\mathcal S}_{d+1}(BG)\to\R/\Z,
  \tag{12.90}
\end{equation}
或者等价地
\begin{equation}
  I_T\in\Hom\bigl(\Omega^{\mathcal S}_{d+1}(BG),\R/\Z\bigr).
  \tag{12.91}
\end{equation}
注意与式 (12.84) 相比，$d+2$ 变成了 $d+1$，$\Z$ 变成了 $\R/\Z$。

当 $\Omega^{\mathcal S}_{d+1}(BG)\simeq\Z^k\oplus\bigoplus_i\Z_{n_i}$ 时，有
\begin{equation}
  \Hom\bigl(\Omega^{\mathcal S}_{d+1}(BG),\R/\Z\bigr)
  \simeq(\R/\Z)^k\oplus\bigoplus_i\Z_{n_i}.
  \tag{12.92}
\end{equation}
其中 $(\R/\Z)^k$ 这一部分可以连续地改变，因此在按连续形变分类相互作用可逆相时并不重要。重要的是 $\bigoplus_i\Z_{n_i}$ 这一部分，即挠部分。对阿贝尔群 $A$，$\Hom(A,\R/\Z)$ 的挠部分也称为 $\operatorname{Ext}(A,\Z)$。所以 $J_T$ 为零的可逆相的形变类确定一个类
\begin{equation}
  [I_T]\in\operatorname{Ext}(\Omega^{\mathcal S}_{d+1}(BG),\Z)
  =\Hom(\Omega^{\mathcal S}_{d+1}(BG),\R/\Z)\ \text{的挠部分}.
  \tag{12.93}
\end{equation}

把时空结构为 $\mathcal S$、对称性为 $G$ 的相互作用可逆相的形变类构成的群记作 $\operatorname{Inv}^{d+1}_{\mathcal S}(BG)$。到此为止的论证告诉我们，这个群位于短正合列
\begin{equation}
  0\to\operatorname{Ext}(\Omega^{\mathcal S}_{d+1}(BG),\Z)
  \to\operatorname{Inv}^{d+1}_{\mathcal S}(BG)
  \xrightarrow{J}\Hom(\Omega^{\mathcal S}_{d+2}(BG),\Z)\to0
  \tag{12.94}
\end{equation}
之中。

% [原文 p.188]
\begin{center}
  \small
  \setlength{\tabcolsep}{4pt}
  \begin{tabular}{c||cc|ccc|ccc}
    \toprule
    $d$
    & $\mho^D_{\mathrm{unoriented}}$ & $\mho^d_{\mathrm{oriented}}$
    & $\mho^d_{\mathrm{pin}^-}$ & $\mho^d_{\mathrm{spin}}$ & $\mho^d_{\mathrm{pin}^+}$
    & $\mho^d_{\mathrm{spin}}(B\Z_2)$ & $\mho^d_{\mathrm{spin}}(BU(1))$
    & $\mho^d_{\mathrm{spin}\text{-}c}$ \\
    \midrule
    0 & $0$ & $\Z$ & $0$ & $\Z$ & $0$ & $\Z$ & $\Z$ & $\Z$ \\
    1 & $\Z_2$ & $0$ & $\Z_2$ & $0$ & $\Z_2$ & $0$ & $0$ & $0$ \\
    2 & $0$ & $0$ & $\Z_2$ & $\Z_2$ & $0$ & $\Z_2^2$ & $\Z\oplus\Z_2$ & $\Z$ \\
    3 & $\Z_2$ & $0$ & $\Z_8$ & $\Z_2$ & $\Z_2$ & $\Z_2^2$ & $\Z_2$ & $0$ \\
    4 & $0$ & $\Z$ & $0$ & $\Z$ & $\Z_2$ & $\Z\oplus\Z_8$ & $\Z^2$ & $\Z^2$ \\
    5 & $\Z_2^2$ & $0$ & $0$ & $0$ & $\Z_{16}$ & $0$ & $0$ & $0$ \\
    6 & $\Z_2$ & $\Z_2$ & $0$ & $0$ & $0$ & $0$ & $\Z^2$ & $\Z^2$ \\
    \bottomrule
  \end{tabular}

  \smallskip
  {\small 表 11：表 10 中所给配边群的 Anderson 对偶。为节省空间，我们把 $(I_\Z\Omega^{\mathcal S})^d$ 记作 $\mho^d_{\mathcal S}$。\par}
\end{center}

这个短正合列也许看着吓人，但在抽象群的层面上，它不过意味着
\begin{equation}
  \operatorname{Inv}^{d+1}_{\mathcal S}(BG)
  =\bigl(\Omega^{\mathcal S}_{d+1}(BG)\ \text{的挠部分}\bigr)
  \oplus\bigl(\Omega^{\mathcal S}_{d+2}(BG)\ \text{的非挠部分}\bigr).
  \tag{12.95}
\end{equation}

其实，数学上有一个概念叫配边群的 Anderson 对偶，记作 $(I_\Z\Omega^{\mathcal S})^{d+2}(BG)$，它位于同一个正合列中：
\begin{equation}
  0\to\operatorname{Ext}(\Omega^{\mathcal S}_{d+1}(BG),\Z)
  \to(I_\Z\Omega^{\mathcal S})^{d+2}(BG)
  \xrightarrow{J}\Hom(\Omega^{\mathcal S}_{d+2}(BG),\Z)\to0,
  \tag{12.96}
\end{equation}
这把我们引向如下结论（或猜想）：
\begin{equation}
  \operatorname{Inv}^{d+1}_{\mathcal S}(BG)=(I_\Z\Omega^{\mathcal S})^{d+2}(BG).
  \tag{12.97}
\end{equation}
表 10 中所给各配边群的 Anderson 对偶列在表 11 中。

我们总应当能把自由费米子相看作相互作用费米子相。这应当给出联系自由分类与相互作用分类的同态。下面讨论三种情形：

\begin{itemize}
  \item 对称类型 $i=1$ 即 BDI，对应 pin$^-$。在 $d+1$ 维，这应当给出映射
  \begin{equation}
    KO^{d-1}(\pt)\to(I_\Z\Omega^{\mathrm{pin}^-})^{d+2}(\pt).
    \tag{12.98}
  \end{equation}
  \item 对称类型 $i=2$ 即 D，对应 spin。在 $d+1$ 维，这给出
  \begin{equation}
    KO^{d-2}(\pt)\to(I_\Z\Omega^{\mathrm{spin}})^{d+2}(\pt).
    \tag{12.99}
  \end{equation}
  \item 对称类型 $i=3$ 即 DIII，对应 pin$^+$。在 $d+1$ 维，这给出
  \begin{equation}
    KO^{d-3}(\pt)\to(I_\Z\Omega^{\mathrm{pin}^+})^{d+2}(\pt).
    \tag{12.100}
  \end{equation}
\end{itemize}

把表 9 的三行与表 11 的三列作比较。要想更充分地理解这些同态，我们需要多学一些代数拓扑。

\takeaway{整值的 $J_T$ 只看到配边群的自由部分；当 $J_T=0$ 时 $I_T$ 本身是取值 $\R/\Z$ 的配边不变量，其形变类落在 $0\to\operatorname{Ext}(\Omega_{d+1},\Z)\to\operatorname{Inv}^{d+1}\to\Hom(\Omega_{d+2},\Z)\to0$ 中，于是相互作用可逆相的分类归结为配边群的 Anderson 对偶。}

\chapterreview{本章先把 eta 不变量定义为质量无穷大的自由费米子配分函数的相位部分，即谱中符号函数正规化和的一半，它在模 $1$ 意义下良定义；APS 指标定理把它写成体上的指标加示性类的积分。再由三条性质公设出可逆相低能有效作用量的相位部分：不交并相加、定向反转取负、在边界上等于示性类 $J_T$ 的积分。引进配边群 $\Omega^{\mathcal S}_d(X)$ 后，整值的 $J_T$ 给出同态 $\Omega^{\mathcal S}_{d+2}(BG)\to\Z$，它只看到自由部分；而当 $J_T$ 为零时，$I_T$ 本身是取值在 $\R/\Z$ 的配边不变量，其形变类落在短正合列 $0\to\operatorname{Ext}(\Omega^{\mathcal S}_{d+1},\Z)\to\operatorname{Inv}^{d+1}_{\mathcal S}\to\Hom(\Omega^{\mathcal S}_{d+2},\Z)\to0$ 中。于是相互作用可逆相的分类归结为配边群的 Anderson 对偶，表 10、表 11 给出了具体的群。}




